\begin{paracol}{2}
\section*{QA3a. Theoretical Q\&A --- Modalities \& Reconstruction}

\textbf{Q1. What types of input data exist in medical imaging? List and briefly describe each.}\\
Medical imaging produces several data types depending on modality. \textbf{2D projection images} (X-ray radiographs) are single shadow images formed by attenuating X-rays. \textbf{Tomographic slice stacks} (CT, MRI) are 2D cross-sections that combine into a 3D \textbf{volumetric scalar field} $f(x,y,z)$ stored as voxels with intensity values. \textbf{Sinograms} are raw projection data $p(\rho,\theta)$ from CT before reconstruction. \textbf{k-space data} are raw frequency-domain samples acquired by MRI before inverse Fourier transform. \textbf{Time-series volumes} (fMRI, 4D-CT) add a temporal axis $f(x,y,z,t)$. \textbf{Functional/molecular data} (PET, SPECT) encode tracer concentration. Finally, \textbf{point clouds and surface meshes} arise after segmentation for surgical planning. All are typically packaged in the \textbf{DICOM} format with metadata.

\textbf{Q2. Explain how MRI works (physics $\to$ image): magnetic field, RF pulse, T1/T2 relaxation.}\\
MRI exploits the magnetic properties of \textbf{hydrogen protons} in water. (1) A strong static field $B_0$ (1.5--3T) aligns proton spins along its axis, producing net magnetization $M_z$. (2) The protons precess at the \textbf{Larmor frequency} $\omega = \gamma B_0$, where $\gamma$ is the gyromagnetic ratio (42.58 MHz/T for $^1$H). (3) An \textbf{RF pulse} at $\omega$ tips spins into the transverse plane, creating $M_{xy}$ that induces a signal in receiver coils. (4) After the pulse, two relaxations occur: \textbf{T1} (spin-lattice, $M_z$ recovery) and \textbf{T2} (spin-spin, $M_{xy}$ decay). (5) \textbf{Gradient coils} encode spatial position via frequency and phase, filling \textbf{k-space}. (6) An \textbf{inverse 2D Fourier transform} of k-space yields the image. Different tissues have different T1/T2, producing contrast.

\textbf{Q3. What is the difference between T1, T2, FLAIR, and DWI MRI?}\\
\textbf{T1-weighted}: short TR/TE; fat appears bright, water dark; great for anatomy and post-contrast (gadolinium) imaging. \textbf{T2-weighted}: long TR/TE; water/edema appears bright, ideal for detecting pathology, inflammation, tumors. \textbf{FLAIR} (Fluid-Attenuated Inversion Recovery) is a T2 variant that nulls free CSF signal so periventricular lesions (e.g., MS plaques) stand out. \textbf{DWI} (Diffusion-Weighted Imaging) measures Brownian motion of water molecules using bipolar gradients; restricted diffusion (acute stroke, dense tumors) appears bright. Each sequence highlights different tissue properties from the same hardware by varying \textbf{TR, TE, TI, and gradient schemes}.

\textbf{Q4. Explain how CT scanner works (X-ray, sinogram, projections).}\\
A CT scanner has an \textbf{X-ray tube} and a ring of detectors that rotate around the patient. (1) X-rays pass through the body; tissues attenuate them following Beer's law $I = I_0 e^{-\int \mu(s)\,ds}$. (2) Each detector measures the line integral $p_\theta(\rho) = \int \mu(x,y)\,d\ell$ along that ray. (3) Rotating from $\theta=0$ to $\pi$ collects many projections forming a \textbf{sinogram} (a point in space traces a sinusoid). (4) Reconstruction (FBP or iterative) inverts the Radon transform to recover the 2D map $\mu(x,y)$. (5) Values are normalized to \textbf{Hounsfield Units}. Modern \textbf{helical multi-detector CT} acquires whole volumes in seconds.

\textbf{Q5. What are Hounsfield Units? Give typical values for air, water, soft tissue, bone.}\\
\textbf{Hounsfield Units (HU)} normalize CT attenuation relative to water:
$$\mathrm{HU} = 1000 \cdot \frac{\mu - \mu_w}{\mu_w}$$
By definition, \textbf{water = 0 HU} and \textbf{air = $-1000$ HU}. Typical values: \textbf{fat} $-100$ to $-50$, \textbf{soft tissue} $+30$ to $+80$, \textbf{blood} $+30$ to $+45$, \textbf{muscle} $+10$ to $+40$, \textbf{cancellous bone} $+300$ to $+400$, \textbf{cortical bone} $+700$ to $+3000$, \textbf{metal/contrast} $> +1000$. HU windowing (window/level) maps a narrow HU range to grayscale: e.g., bone window $W=2000, L=400$; lung window $W=1500, L=-600$. This standardization lets clinicians compare scans across machines.

\textbf{Q6. Compare MRI vs CT --- strengths, weaknesses.}\\
\textbf{CT strengths}: fast (seconds), excellent for \textbf{bone, lung, acute hemorrhage, trauma}, widely available, cheap. \textbf{CT weaknesses}: \textbf{ionizing radiation}, poor soft-tissue contrast, beam-hardening artifacts near metal. \textbf{MRI strengths}: \textbf{superb soft-tissue contrast} (brain, spine, joints, pelvis), no ionizing radiation, multi-parametric (T1/T2/DWI/fMRI), arbitrary slice planes. \textbf{MRI weaknesses}: slow (minutes), expensive, contraindicated for pacemakers/ferromagnetic implants, claustrophobia, sensitive to motion. \textbf{Use CT} for emergency, bone, lung, and quick screening; \textbf{use MRI} for neurological, musculoskeletal, oncologic, and cardiac soft-tissue evaluation.

\textbf{Q7. How does ultrasound (US) work? When is it preferred?}\\
A piezoelectric transducer emits high-frequency \textbf{sound pulses (2--15 MHz)} into tissue and receives reflected echoes. The time-of-flight gives depth ($d = ct/2$, $c \approx 1540$ m/s), and echo amplitude maps to brightness in \textbf{B-mode}. \textbf{Doppler mode} measures blood-flow velocity from frequency shifts. Ultrasound is preferred for \textbf{obstetrics} (no radiation safe for fetus), \textbf{cardiac echo}, \textbf{vascular flow}, \textbf{abdominal organs}, and \textbf{point-of-care} bedside imaging. Strengths: real-time, portable, cheap, no radiation. Weaknesses: poor through bone/air, operator-dependent, limited field of view.

\textbf{Q8. How does PET imaging work? What does FDG measure?}\\
\textbf{Positron Emission Tomography} injects a radioactive tracer that emits \textbf{positrons}; each positron annihilates with an electron producing two \textbf{511 keV gamma photons} traveling 180$^\circ$ apart. A ring of detectors records \textbf{coincidence events}, and reconstruction (OSEM iterative) yields a 3D map of tracer concentration. \textbf{FDG ($^{18}$F-fluorodeoxyglucose)} is a glucose analog trapped in metabolically active cells, revealing tumors, inflammation, and brain activity. PET shows \textbf{function/metabolism}, not anatomy. Other tracers (PSMA, amyloid, $^{68}$Ga-DOTATATE) target specific receptors. Spatial resolution is modest ($\sim$4--6 mm) but molecular sensitivity is unmatched.

\textbf{Q9. PET-CT vs PET-MRI?}\\
\textbf{PET-CT} fuses functional PET with anatomic CT acquired sequentially in the same gantry; CT also provides \textbf{attenuation correction}. It is fast, robust, and standard for oncology staging, but adds CT radiation dose. \textbf{PET-MRI} replaces CT with MRI, giving \textbf{superior soft-tissue contrast} and no extra radiation, ideal for pediatric, neurological, and pelvic cancers. PET-MRI is more expensive, slower, and attenuation correction is harder (MRI doesn't directly measure electron density --- requires Dixon-based or atlas methods). Choice depends on indication, dose concern, and availability.

\textbf{Q10. What is DICOM? Why is it the standard?}\\
\textbf{DICOM} (Digital Imaging and Communications in Medicine) is the universal standard for storing and transmitting medical images. Each \textbf{DICOM file} contains pixel data plus a header of \textbf{tagged metadata}: patient ID, study/series UIDs, modality, acquisition parameters (TR, TE, kVp, mAs), pixel spacing, slice thickness, orientation matrix. It defines \textbf{network services} (C-STORE, C-FIND, C-MOVE) for PACS communication. Standardization ensures \textbf{interoperability} across vendors, enables longitudinal comparison, supports AI pipelines, and provides legal/audit trails. Without DICOM, every scanner would speak its own dialect.

\textbf{Q11. What is fMRI and what does it measure?}\\
\textbf{Functional MRI} measures brain activity via the \textbf{BOLD signal} (Blood-Oxygen-Level-Dependent). When neurons fire, local blood flow increases beyond oxygen consumption, raising the ratio of oxyhemoglobin (diamagnetic) to deoxyhemoglobin (paramagnetic) and lengthening T2*. A \textbf{T2*-weighted EPI sequence} captures this small ($\sim$1--5\%) signal change every 1--2 seconds. Subjects perform tasks (block or event-related design) and statistical models (GLM) localize activated regions. fMRI provides \textbf{spatial resolution} of millimeters but \textbf{temporal resolution} of seconds (vs EEG's ms). Used for presurgical planning, cognitive neuroscience, and resting-state connectivity.

\textbf{Q12. What is SPECT? How does it differ from PET?}\\
\textbf{SPECT} (Single-Photon Emission CT) uses tracers emitting a \textbf{single gamma photon} (e.g., $^{99m}$Tc, $^{123}$I) detected by a rotating gamma camera with a \textbf{collimator}. Reconstruction is similar to CT (FBP/OSEM) on projection data. Differences vs PET: SPECT photons are single (not coincidence pairs), so \textbf{collimators} are needed, reducing sensitivity by $\sim$100$\times$. Resolution is lower ($\sim$10 mm vs 4 mm), but tracers are cheaper and longer-lived (no cyclotron needed). SPECT is widely used for \textbf{cardiac perfusion, bone scans, thyroid, brain (DAT-scan)}.

\switchcolumn

\section*{QA3a. Câu hỏi lý thuyết --- Phương thức tạo ảnh \& Tái tạo}

\textbf{Câu 1. Có những loại dữ liệu đầu vào nào trong y học hình ảnh?}\\
Y học hình ảnh sinh ra nhiều loại dữ liệu tùy theo modality. \textbf{Ảnh chiếu 2D} (X-quang) là một ảnh bóng do tia X bị suy giảm. \textbf{Chồng lát cắt tomographic} (CT, MRI) là các mặt cắt 2D ghép thành \textbf{trường vô hướng 3D} $f(x,y,z)$ lưu dưới dạng voxel. \textbf{Sinogram} là dữ liệu chiếu thô $p(\rho,\theta)$ từ CT trước khi tái tạo. \textbf{k-space} là mẫu thô miền tần số do MRI thu thập trước khi biến đổi Fourier ngược. \textbf{Khối thời gian} (fMRI, 4D-CT) thêm trục thời gian $f(x,y,z,t)$. \textbf{Dữ liệu chức năng/phân tử} (PET, SPECT) mã hóa nồng độ chất đánh dấu. Cuối cùng, \textbf{point cloud và mesh bề mặt} xuất hiện sau segmentation phục vụ phẫu thuật. Tất cả đóng gói trong định dạng \textbf{DICOM} kèm metadata.

\textbf{Câu 2. MRI hoạt động thế nào (vật lý $\to$ ảnh)?}\\
MRI khai thác từ tính của \textbf{proton hydrogen} trong nước. (1) Từ trường tĩnh mạnh $B_0$ (1.5--3T) căn chỉnh spin proton dọc trục, tạo từ hóa $M_z$. (2) Proton tiến động (precess) ở \textbf{tần số Larmor} $\omega = \gamma B_0$ (42.58 MHz/T với $^1$H). (3) Một \textbf{xung RF} ở $\omega$ lật spin sang mặt phẳng ngang, tạo $M_{xy}$ cảm ứng tín hiệu trong cuộn thu. (4) Sau xung, hai quá trình hồi phục: \textbf{T1} (spin-mạng, hồi $M_z$) và \textbf{T2} (spin-spin, suy giảm $M_{xy}$). (5) \textbf{Cuộn gradient} mã hóa vị trí qua tần số và pha, lấp đầy \textbf{k-space}. (6) Biến đổi \textbf{Fourier ngược 2D} của k-space cho ảnh. Mô khác nhau có T1/T2 khác nhau, tạo độ tương phản.

\textbf{Câu 3. Khác biệt T1, T2, FLAIR, DWI?}\\
\textbf{T1-weighted}: TR/TE ngắn; mỡ sáng, nước tối; tốt cho giải phẫu và sau tiêm gadolinium. \textbf{T2-weighted}: TR/TE dài; nước/phù sáng, lý tưởng để phát hiện bệnh lý, viêm, u. \textbf{FLAIR} là biến thể T2 triệt tín hiệu CSF tự do để tổn thương cạnh não thất (ví dụ MS) nổi bật. \textbf{DWI} đo chuyển động Brown của phân tử nước qua gradient lưỡng cực; khuếch tán hạn chế (đột quỵ cấp, u đặc) sáng. Mỗi sequence làm nổi bật tính chất mô khác nhau từ cùng phần cứng bằng cách thay đổi \textbf{TR, TE, TI, và gradient}.

\textbf{Câu 4. CT scanner hoạt động thế nào?}\\
Máy CT có \textbf{ống tia X} và vành detector quay quanh bệnh nhân. (1) Tia X xuyên cơ thể; mô làm suy giảm theo định luật Beer $I = I_0 e^{-\int \mu(s)\,ds}$. (2) Mỗi detector đo tích phân đường $p_\theta(\rho) = \int \mu(x,y)\,d\ell$. (3) Quay từ $\theta=0$ đến $\pi$ thu nhiều projection tạo \textbf{sinogram} (một điểm vẽ thành đường sin). (4) Tái tạo (FBP hoặc iterative) đảo ngược biến đổi Radon để khôi phục bản đồ 2D $\mu(x,y)$. (5) Giá trị chuẩn hóa thành \textbf{Hounsfield Units}. CT đa dãy xoắn ốc hiện đại thu cả khối trong vài giây.

\textbf{Câu 5. Hounsfield Units là gì? Giá trị điển hình?}\\
\textbf{Hounsfield Units (HU)} chuẩn hóa độ suy giảm CT theo nước:
$$\mathrm{HU} = 1000 \cdot \frac{\mu - \mu_w}{\mu_w}$$
Theo định nghĩa, \textbf{nước = 0 HU} và \textbf{khí = $-1000$ HU}. Giá trị điển hình: \textbf{mỡ} $-100$ đến $-50$, \textbf{mô mềm} $+30$ đến $+80$, \textbf{máu} $+30$ đến $+45$, \textbf{cơ} $+10$ đến $+40$, \textbf{xương xốp} $+300$ đến $+400$, \textbf{xương vỏ} $+700$ đến $+3000$, \textbf{kim loại/cản quang} $> +1000$. Windowing HU (window/level) ánh xạ dải HU hẹp sang thang xám: ví dụ window xương $W=2000, L=400$; window phổi $W=1500, L=-600$. Chuẩn hóa này cho phép so sánh giữa các máy.

\textbf{Câu 6. So sánh MRI vs CT.}\\
\textbf{Ưu điểm CT}: nhanh (vài giây), tuyệt vời cho \textbf{xương, phổi, xuất huyết cấp, chấn thương}, phổ biến, rẻ. \textbf{Nhược điểm CT}: \textbf{tia ion hóa}, độ tương phản mô mềm kém, nhiễu beam-hardening gần kim loại. \textbf{Ưu điểm MRI}: \textbf{độ tương phản mô mềm xuất sắc} (não, cột sống, khớp, vùng chậu), không tia ion hóa, đa tham số (T1/T2/DWI/fMRI), mặt cắt tùy ý. \textbf{Nhược điểm MRI}: chậm (vài phút), đắt, chống chỉ định với máy tạo nhịp/cấy ghép sắt từ, gây sợ hãi không gian kín, nhạy với chuyển động. \textbf{Dùng CT} cho cấp cứu, xương, phổi, sàng lọc nhanh; \textbf{dùng MRI} cho thần kinh, cơ-xương-khớp, ung bướu mô mềm, tim.

\textbf{Câu 7. Siêu âm hoạt động thế nào? Khi nào dùng?}\\
Đầu dò áp điện phát \textbf{xung âm tần cao (2--15 MHz)} vào mô và nhận sóng dội. Thời gian bay cho độ sâu ($d = ct/2$, $c \approx 1540$ m/s), biên độ dội ánh xạ thành độ sáng \textbf{B-mode}. \textbf{Doppler} đo vận tốc dòng máu qua dịch tần số. Siêu âm được ưu tiên cho \textbf{sản khoa} (an toàn cho thai), \textbf{siêu âm tim}, \textbf{mạch máu}, \textbf{tạng bụng}, và \textbf{cấp cứu tại giường}. Ưu: thời gian thực, di động, rẻ, không tia. Nhược: kém qua xương/khí, phụ thuộc kỹ thuật viên, trường nhìn hẹp.

\textbf{Câu 8. PET hoạt động thế nào? FDG đo gì?}\\
\textbf{PET (Positron Emission Tomography)} tiêm chất đánh dấu phóng xạ phát \textbf{positron}; mỗi positron hủy với electron sinh hai \textbf{photon gamma 511 keV} bay ngược chiều 180$^\circ$. Vành detector ghi \textbf{sự kiện trùng phùng} (coincidence), và tái tạo (OSEM iterative) cho bản đồ 3D nồng độ chất đánh dấu. \textbf{FDG ($^{18}$F-fluorodeoxyglucose)} là chất tương tự glucose bị giữ trong tế bào chuyển hóa cao, chỉ điểm u, viêm, hoạt động não. PET cho thấy \textbf{chức năng/chuyển hóa}, không phải giải phẫu. Các tracer khác (PSMA, amyloid, $^{68}$Ga-DOTATATE) nhắm receptor cụ thể. Độ phân giải không gian khiêm tốn ($\sim$4--6 mm) nhưng độ nhạy phân tử vô địch.

\textbf{Câu 9. PET-CT vs PET-MRI?}\\
\textbf{PET-CT} kết hợp PET chức năng với CT giải phẫu thu liên tiếp trong cùng gantry; CT cũng cung cấp \textbf{hiệu chỉnh suy giảm}. Nhanh, mạnh mẽ, chuẩn cho phân giai đoạn ung thư, nhưng thêm liều xạ CT. \textbf{PET-MRI} thay CT bằng MRI, cho \textbf{tương phản mô mềm vượt trội} và không thêm tia, lý tưởng cho ung thư nhi, thần kinh, vùng chậu. PET-MRI đắt hơn, chậm hơn, hiệu chỉnh suy giảm khó hơn (MRI không đo trực tiếp mật độ điện tử --- cần phương pháp Dixon hoặc atlas). Lựa chọn tùy chỉ định, lo ngại liều, và sẵn có.

\textbf{Câu 10. DICOM là gì? Tại sao là chuẩn?}\\
\textbf{DICOM} (Digital Imaging and Communications in Medicine) là chuẩn toàn cầu để lưu và truyền ảnh y tế. Mỗi \textbf{file DICOM} chứa dữ liệu pixel cộng header \textbf{metadata gắn tag}: ID bệnh nhân, UID study/series, modality, tham số thu (TR, TE, kVp, mAs), pixel spacing, slice thickness, ma trận hướng. Định nghĩa \textbf{dịch vụ mạng} (C-STORE, C-FIND, C-MOVE) cho giao tiếp PACS. Chuẩn hóa đảm bảo \textbf{liên thông} giữa các hãng, cho phép so sánh dài hạn, hỗ trợ pipeline AI, và tạo dấu vết pháp lý/kiểm toán. Không có DICOM, mỗi máy sẽ nói tiếng riêng.

\textbf{Câu 11. fMRI là gì và đo gì?}\\
\textbf{fMRI} đo hoạt động não qua \textbf{tín hiệu BOLD} (Blood-Oxygen-Level-Dependent). Khi neuron phóng điện, lưu lượng máu cục bộ tăng vượt mức tiêu thụ oxy, làm tăng tỉ lệ oxyhemoglobin (nghịch từ) so với deoxyhemoglobin (thuận từ) và kéo dài T2*. \textbf{Sequence EPI T2*-weighted} bắt thay đổi tín hiệu nhỏ ($\sim$1--5\%) mỗi 1--2 giây. Đối tượng thực hiện task (block hoặc event-related) và mô hình thống kê (GLM) định vị vùng kích hoạt. fMRI cho \textbf{độ phân giải không gian} milimét nhưng \textbf{độ phân giải thời gian} chỉ vài giây (so với EEG mili giây). Dùng cho lập kế hoạch tiền phẫu, khoa học thần kinh nhận thức, và kết nối resting-state.

\textbf{Câu 12. SPECT là gì? Khác PET ra sao?}\\
\textbf{SPECT} (Single-Photon Emission CT) dùng chất đánh dấu phát \textbf{một photon gamma đơn} (ví dụ $^{99m}$Tc, $^{123}$I) được phát hiện bởi gamma camera xoay với \textbf{collimator}. Tái tạo tương tự CT (FBP/OSEM) trên dữ liệu projection. Khác biệt với PET: photon SPECT là đơn (không phải cặp trùng phùng), nên cần \textbf{collimator}, giảm độ nhạy $\sim$100 lần. Độ phân giải thấp hơn ($\sim$10 mm vs 4 mm), nhưng tracer rẻ hơn và sống lâu hơn (không cần cyclotron). SPECT dùng nhiều cho \textbf{tưới máu cơ tim, xạ hình xương, tuyến giáp, não (DAT-scan)}.

\switchcolumn*

\textbf{Q13. Reconstruct a 2D CT slice via Filtered Back Projection (FBP).}\\
\textbf{FBP} inverts the Radon transform. Steps: (1) \textbf{Acquire} projections $p_\theta(\rho)$ for $\theta \in [0,\pi]$. (2) \textbf{1D Fourier transform} each projection: $P_\theta(\nu) = \mathcal{F}\{p_\theta\}$. (3) \textbf{Apply ramp filter} $|\nu|$ (often windowed by Ram-Lak, Shepp-Logan, or Hann to reduce noise): $\tilde{P}_\theta(\nu) = |\nu| P_\theta(\nu)$. (4) \textbf{Inverse 1D FT} to get filtered projections $\tilde{p}_\theta(\rho)$. (5) \textbf{Back-project} by smearing each filtered projection along its ray and integrating over angles:
$$f(x,y) = \int_0^\pi \tilde{p}_\theta(x\cos\theta + y\sin\theta)\,d\theta$$
The ramp filter compensates for the $1/|\nu|$ blurring inherent in plain back-projection. FBP is fast, analytic, but noisy at low dose.

\textbf{Q14. What is the Radon transform / sinogram?}\\
The \textbf{Radon transform} maps a 2D function $f(x,y)$ to its line integrals:
$$p_\theta(\rho) = \int_{-\infty}^{\infty} f(\rho\cos\theta - s\sin\theta,\,\rho\sin\theta + s\cos\theta)\,ds$$
where $\rho$ is the perpendicular distance and $\theta$ the angle of the line. Stacking $p_\theta(\rho)$ as a 2D image with $\rho$ on the horizontal and $\theta$ on the vertical axis gives the \textbf{sinogram}. A point object traces a sinusoid (hence the name). CT reconstruction = inverting this transform. The Radon transform underlies all tomography and is invertible only if projections cover at least 180$^\circ$ (parallel beam).

\textbf{Q15. What is the Central Slice Theorem?}\\
The \textbf{Central Slice Theorem} (Fourier Slice Theorem) states: the \textbf{1D Fourier transform of a projection at angle $\theta$} equals a \textbf{1D radial slice through the 2D Fourier transform of the object} at the same angle:
$$\mathcal{F}_1\{p_\theta(\rho)\}(\nu) = F(\nu\cos\theta,\,\nu\sin\theta)$$
This is the theoretical foundation of CT reconstruction: collecting projections at many angles fills the 2D Fourier plane radially, and an inverse 2D FT recovers $f$. Because radial sampling is denser near the origin, a $|\nu|$ \textbf{ramp filter} compensates for this density before back-projection --- which is exactly FBP.

\textbf{Q16. How is MRI reconstructed from k-space?}\\
\textbf{k-space} is the 2D (or 3D) frequency domain of the image. Each MRI readout fills one trajectory in k-space (line for Cartesian, spiral, radial). Steps: (1) \textbf{Acquire} all required k-space samples by sweeping phase-encoding gradients and frequency-encoding readouts. (2) For Cartesian sampling, apply \textbf{inverse 2D FFT}: $I(x,y) = \mathcal{F}^{-1}\{S(k_x,k_y)\}$. (3) For non-Cartesian (radial, spiral), \textbf{regrid} to Cartesian then FFT, or use \textbf{NUFFT}. (4) Combine multi-coil data using \textbf{sum-of-squares} or sensitivity-weighted combination. (5) Apply intensity correction. Center of k-space encodes contrast/low frequencies; periphery encodes edges/detail. Undersampling enables acceleration (SENSE, GRAPPA, compressed sensing).

\textbf{Q17. Build a 3D volume from a stack of 2D MRI slices.}\\
(1) \textbf{Read DICOM} stack, sort by \texttt{InstanceNumber} or \texttt{SliceLocation}. (2) Use header metadata (\texttt{PixelSpacing}, \texttt{SliceThickness}, \texttt{ImagePositionPatient}, \texttt{ImageOrientationPatient}) to compute world-space voxel positions. (3) \textbf{Resample} to isotropic spacing if needed (trilinear). (4) Store as a 3D array $V[i,j,k]$. (5) \textbf{Visualize} via: (a) \textbf{Multi-planar reformat (MPR)} --- axial/coronal/sagittal slices; (b) \textbf{Volume rendering} with a transfer function; (c) \textbf{Marching Cubes} surface extraction at an iso-value; (d) \textbf{MIP} (Maximum Intensity Projection) for vascular display. Tools: VTK, ITK-SNAP, 3D Slicer.

\textbf{Q18. Marching Cubes: 256 cases, 15 unique.}\\
\textbf{Marching Cubes} (Lorensen \& Cline 1987) extracts an iso-surface $f(x,y,z) = c$ from a scalar volume. Steps: (1) Slide a logical cube over each set of 8 neighboring voxels. (2) For each corner, mark inside ($f \geq c$) or outside ($f < c$) --- 8 bits. (3) The 8-bit code (\textbf{$2^8 = 256$ cases}) indexes a precomputed \textbf{triangulation lookup table}. (4) For each edge crossed by the surface, place a vertex by \textbf{linear interpolation}:
$$\mathbf{v} = \mathbf{p}_1 + \frac{c - f_1}{f_2 - f_1}(\mathbf{p}_2 - \mathbf{p}_1)$$
(5) Connect vertices into triangles per the lookup table; advance to next cube. By symmetries (rotations + complement), 256 cases reduce to \textbf{15 unique topologies}. Output: a triangle mesh approximating the iso-surface.

\textbf{Q19. Volume rendering (DVR) vs Surface rendering (Marching Cubes).}\\
\textbf{Direct Volume Rendering} ray-casts through the volume, accumulating color/opacity via a \textbf{transfer function}: $C_{out} = \sum_i C_i \alpha_i \prod_{j<i}(1-\alpha_j)$. It shows \textbf{all internal structure semi-transparently}, preserves all data, ideal for soft tissue and exploration, but slow and harder to measure. \textbf{Surface rendering (Marching Cubes)} extracts a polygon mesh at one iso-value; fast to render afterwards on GPU, easy to measure/3D-print/cut, but \textbf{commits to one threshold} and discards data. \textbf{Use DVR} for diagnostic exploration, multi-tissue display; \textbf{use surface} for surgical planning, 3D printing, simulation, mesh-based analysis.

\textbf{Q20. What is a transfer function in volume rendering?}\\
A \textbf{transfer function} maps scalar voxel values (and optionally gradient magnitude) to \textbf{color and opacity}: $\mathrm{TF}(v) \to (R,G,B,\alpha)$. It controls which tissues are visible: e.g., assign $\alpha=0$ to soft tissue and $\alpha=0.8$ red to bone in a CT to show only the skeleton. \textbf{1D TFs} use intensity only; \textbf{2D TFs} add gradient magnitude to highlight boundaries. Designing a good TF is the \textbf{hardest part of DVR} --- often interactive with histogram widgets. Preset TFs for CT (bone, lung, vessel) and MRI ship with viewers. The TF is essentially a visual segmentation without explicit labels.

\textbf{Q21. Iterative reconstruction (MLEM, OSEM) vs FBP.}\\
\textbf{FBP} is fast, analytic, exact for noiseless full data but \textbf{amplifies noise} (ramp filter) and produces streaks at low dose / sparse views. \textbf{MLEM} (Maximum Likelihood Expectation Maximization) models photon statistics (Poisson) and iteratively updates: $f^{(n+1)} = f^{(n)} \cdot A^T(p / Af^{(n)}) / A^T \mathbf{1}$. \textbf{OSEM} (Ordered Subset EM) accelerates MLEM by partitioning projections into subsets, $\sim$10$\times$ speedup. Iterative methods incorporate \textbf{system models} (geometry, scatter, attenuation) and \textbf{regularization} (TV, edge-preserving), giving better SNR at low dose, but are \textbf{slower} and can over-smooth. Modern CT/PET use iterative as default; deep-learning reconstruction is emerging.

\textbf{Q22. SENSE and GRAPPA parallel imaging for MRI.}\\
Parallel imaging accelerates MRI by undersampling k-space and using \textbf{multi-coil arrays} (8--64 coils), each with a known spatial sensitivity. \textbf{SENSE} (SENSitivity Encoding) works in image space: undersampled k-space gives an aliased image per coil; SENSE solves a linear system with the coil sensitivities $C_i$ to unfold the aliases at each pixel. \textbf{GRAPPA} (GeneRalized Autocalibrating Partial Parallel Acquisition) works in k-space: it uses a fully sampled \textbf{auto-calibration} region to learn linear weights that synthesize missing k-space lines from neighboring acquired lines across coils. Acceleration factors $R = 2$--$4$ are routine, reducing scan time at the cost of $\sqrt{R}$ SNR.

\textbf{Q23. What is compressed sensing in MRI?}\\
\textbf{Compressed sensing (CS)} reconstructs MR images from \textbf{far fewer k-space samples} than Nyquist requires, exploiting two facts: (1) MR images are \textbf{sparse} in some transform domain (wavelets, total variation, finite differences); (2) random/incoherent sampling spreads aliasing as noise-like artifacts. CS solves:
$$\min_x \|\Psi x\|_1 \quad \text{s.t.} \quad \|F_u x - y\|_2 \leq \epsilon$$
where $\Psi$ is the sparsifying transform, $F_u$ undersampled Fourier, $y$ measured k-space. Nonlinear iterative solvers (FISTA, ADMM) recover $x$. CS enables 4--10$\times$ acceleration, especially for dynamic MRI (cardiac, contrast). Combined with parallel imaging and deep learning (variational networks), it underlies modern fast MRI.

\textbf{Q24. VTK pipeline for medical visualization.}\\
\textbf{VTK} (Visualization Toolkit) uses a data-flow pipeline: \textbf{Source $\to$ Filter $\to$ Mapper $\to$ Actor $\to$ Renderer $\to$ RenderWindow}. For CT: (1) \texttt{vtkDICOMImageReader} loads volume. (2) Filters preprocess: \texttt{vtkImageThreshold}, \texttt{vtkImageGaussianSmooth}. (3) \texttt{vtkMarchingCubes} or \texttt{vtkContourFilter} extracts iso-surface; \textit{or} \texttt{vtkSmartVolumeMapper} does GPU ray-cast DVR. (4) Optionally \texttt{vtkDecimatePro}, \texttt{vtkSmoothPolyDataFilter}. (5) \texttt{vtkPolyDataMapper} converts to graphics primitives. (6) \texttt{vtkActor} adds material/color; \texttt{vtkVolume} for DVR with \texttt{vtkColorTransferFunction} and \texttt{vtkPiecewiseFunction}. (7) \texttt{vtkRenderer} composites; \texttt{vtkRenderWindow} displays; \texttt{vtkRenderWindowInteractor} handles input. 3D Slicer and ParaView are built on VTK.

\textbf{Q25. From CT scan to 3D-printable bone mesh --- pipeline.}\\
(1) \textbf{Acquire} CT DICOM stack. (2) \textbf{Load} into 3D Slicer / ITK-SNAP. (3) \textbf{Window/level} to bone (HU $> +200$). (4) \textbf{Segment} bone by thresholding (HU $> 200$--300) plus \textbf{region-growing} or AI tools to remove vessels with contrast and table. (5) \textbf{Morphological cleanup}: closing, fill holes, island removal. (6) \textbf{Marching Cubes} to generate triangle mesh at iso-value. (7) \textbf{Smooth} (Laplacian or Taubin) and \textbf{decimate} to reasonable triangle count. (8) \textbf{Repair}: ensure manifold, watertight, no self-intersections (MeshLab, Meshmixer). (9) \textbf{Export STL/3MF}. (10) \textbf{Slice} for printing (Cura/PrusaSlicer) and print on FDM/SLA. Used for orthopedic implant planning, education, surgical guides.

\textbf{Q26. How does Marching Cubes help surgical planning?}\\
After CT/MRI segmentation, \textbf{Marching Cubes} converts each labeled organ (tumor, vessel, bone) into a \textbf{closed triangle mesh}. Surgeons can: (1) view the 3D mesh from any angle, rotating in real time on GPU; (2) measure distances, volumes, angles for resection planning; (3) print physical models for tactile rehearsal; (4) feed meshes into \textbf{VR/AR systems} to overlay anatomy on the patient; (5) plan tool trajectories avoiding critical structures. Compared with raw volume display, surface meshes are \textbf{lightweight, easy to manipulate, and 3D-print ready}, making MC the workhorse of patient-specific surgical planning.

\textbf{Q27. How is volume rendering used clinically?}\\
\textbf{Direct Volume Rendering} is used for: (1) \textbf{CT angiography} --- visualize vessels and stenoses in 3D without segmentation; (2) \textbf{Virtual colonoscopy/bronchoscopy} --- fly-through of hollow organs; (3) \textbf{Maxillofacial / orthopedic planning} --- skull and bone overview; (4) \textbf{Radiation therapy planning} --- target and OAR display; (5) \textbf{Patient communication} --- intuitive 3D images for consent. Modern GPU ray-casting renders 512$^3$ volumes at $>$30 fps. Cinematic rendering (path tracing, Monte Carlo) further improves realism. DVR complements MPR slices by giving spatial context without requiring perfect segmentation.

\textbf{Q28. Why is real-time reconstruction important in interventional procedures?}\\
\textbf{Interventional radiology} (catheter navigation, biopsy, ablation, surgery under fluoroscopy or MR) needs imaging \textbf{updated within milliseconds} so the clinician can react. Specifically: (1) \textbf{tool tracking} requires latency $<$ 100 ms; (2) \textbf{respiratory/cardiac motion} demands frame rates $>$ 5 Hz; (3) \textbf{closed-loop control} (MR-guided focused ultrasound, robotic surgery) needs deterministic sub-frame latency. Slow reconstruction means stale guidance, prolonged radiation/anesthesia, and increased complication risk. Hence \textbf{GPU-accelerated FBP, parallel imaging with rapid reconstruction, real-time compressed sensing, and AI-based reconstruction} are critical research areas.

\textbf{Q29. How is image processing used to detect tumor signals in MRI?}\\
A typical \textbf{tumor detection pipeline}: (1) \textbf{Multi-sequence acquisition} (T1, T1+Gd, T2, FLAIR, DWI/ADC) gives complementary contrasts. (2) \textbf{Preprocessing}: bias-field correction (N4), skull stripping, registration of sequences, intensity normalization. (3) \textbf{Feature extraction}: voxel intensities across sequences, texture (GLCM, LBP), gradients, location atlas. (4) \textbf{Segmentation} using thresholding, region growing, level sets, or modern \textbf{CNN/U-Net} models (BraTS challenge). (5) \textbf{Post-processing}: connected components, volume computation, contour overlay. (6) \textbf{Radiomics}: extract shape and texture features for prognosis. Output: tumor mask, volume, and similarity to template lesions, helping radiologists detect, classify, and follow tumors over time.

\switchcolumn

\textbf{Câu 13. Tái tạo lát cắt CT 2D bằng FBP.}\\
\textbf{FBP} đảo ngược biến đổi Radon. Các bước: (1) \textbf{Thu} projection $p_\theta(\rho)$ với $\theta \in [0,\pi]$. (2) \textbf{Fourier 1D} từng projection: $P_\theta(\nu) = \mathcal{F}\{p_\theta\}$. (3) \textbf{Áp ramp filter} $|\nu|$ (thường có cửa sổ Ram-Lak, Shepp-Logan, hoặc Hann để giảm nhiễu): $\tilde{P}_\theta(\nu) = |\nu| P_\theta(\nu)$. (4) \textbf{FT 1D ngược} để được projection đã lọc $\tilde{p}_\theta(\rho)$. (5) \textbf{Back-project} bằng cách trải mỗi projection đã lọc dọc tia và tích phân theo góc:
$$f(x,y) = \int_0^\pi \tilde{p}_\theta(x\cos\theta + y\sin\theta)\,d\theta$$
Ramp filter bù cho làm mờ $1/|\nu|$ vốn có trong back-projection thuần. FBP nhanh, giải tích, nhưng nhiễu khi liều thấp.

\textbf{Câu 14. Biến đổi Radon / sinogram là gì?}\\
\textbf{Biến đổi Radon} ánh xạ hàm 2D $f(x,y)$ thành các tích phân đường:
$$p_\theta(\rho) = \int_{-\infty}^{\infty} f(\rho\cos\theta - s\sin\theta,\,\rho\sin\theta + s\cos\theta)\,ds$$
trong đó $\rho$ là khoảng cách vuông góc và $\theta$ là góc đường thẳng. Xếp $p_\theta(\rho)$ thành ảnh 2D với $\rho$ ngang và $\theta$ dọc cho \textbf{sinogram}. Một điểm vẽ thành đường sin (vì thế tên gọi). Tái tạo CT = đảo ngược biến đổi này. Biến đổi Radon là nền tảng mọi tomography và chỉ khả nghịch nếu projection phủ tối thiểu 180$^\circ$ (chùm song song).

\textbf{Câu 15. Định lý lát cắt trung tâm là gì?}\\
\textbf{Định lý lát cắt trung tâm} (Fourier Slice Theorem) nói: \textbf{biến đổi Fourier 1D của projection ở góc $\theta$} bằng \textbf{lát cắt xuyên tâm 1D qua biến đổi Fourier 2D của vật} ở cùng góc:
$$\mathcal{F}_1\{p_\theta(\rho)\}(\nu) = F(\nu\cos\theta,\,\nu\sin\theta)$$
Đây là nền tảng lý thuyết của tái tạo CT: thu projection ở nhiều góc lấp đầy mặt phẳng Fourier 2D theo tia, và FT 2D ngược khôi phục $f$. Vì lấy mẫu xuyên tâm dày hơn gần gốc, \textbf{ramp filter $|\nu|$} bù mật độ này trước back-projection --- chính là FBP.

\textbf{Câu 16. Tái tạo MRI từ k-space.}\\
\textbf{k-space} là miền tần số 2D (hoặc 3D) của ảnh. Mỗi readout MRI lấp một quỹ đạo trong k-space (dòng cho Cartesian, xoắn ốc, xuyên tâm). Các bước: (1) \textbf{Thu} đủ mẫu k-space bằng quét gradient phase-encoding và readout frequency-encoding. (2) Với mẫu Cartesian, áp \textbf{FFT 2D ngược}: $I(x,y) = \mathcal{F}^{-1}\{S(k_x,k_y)\}$. (3) Với phi Cartesian (xuyên tâm, xoắn ốc), \textbf{regrid} về Cartesian rồi FFT, hoặc dùng \textbf{NUFFT}. (4) Kết hợp dữ liệu nhiều cuộn dùng \textbf{sum-of-squares} hoặc trọng số sensitivity. (5) Hiệu chỉnh cường độ. Tâm k-space mã hóa tương phản/tần thấp; rìa mã hóa cạnh/chi tiết. Lấy mẫu thưa cho phép tăng tốc (SENSE, GRAPPA, compressed sensing).

\textbf{Câu 17. Xây khối 3D từ chồng lát cắt MRI 2D.}\\
(1) \textbf{Đọc DICOM} stack, sắp xếp theo \texttt{InstanceNumber} hoặc \texttt{SliceLocation}. (2) Dùng metadata header (\texttt{PixelSpacing}, \texttt{SliceThickness}, \texttt{ImagePositionPatient}, \texttt{ImageOrientationPatient}) để tính vị trí voxel trong không gian thế giới. (3) \textbf{Resample} sang spacing đẳng hướng nếu cần (trilinear). (4) Lưu thành mảng 3D $V[i,j,k]$. (5) \textbf{Hiển thị} qua: (a) \textbf{Multi-planar reformat (MPR)} --- các lát axial/coronal/sagittal; (b) \textbf{Volume rendering} với transfer function; (c) \textbf{Marching Cubes} trích bề mặt iso-value; (d) \textbf{MIP} (Maximum Intensity Projection) hiển thị mạch máu. Công cụ: VTK, ITK-SNAP, 3D Slicer.

\textbf{Câu 18. Marching Cubes: 256 trường hợp, 15 duy nhất.}\\
\textbf{Marching Cubes} (Lorensen \& Cline 1987) trích iso-surface $f(x,y,z) = c$ từ khối vô hướng. Bước: (1) Trượt khối logic qua mỗi nhóm 8 voxel kề. (2) Mỗi đỉnh đánh dấu trong ($f \geq c$) hoặc ngoài ($f < c$) --- 8 bit. (3) Mã 8-bit (\textbf{$2^8 = 256$ trường hợp}) tra bảng \textbf{lookup tam giác hóa} đã tính trước. (4) Với mỗi cạnh bị bề mặt cắt, đặt đỉnh bằng \textbf{nội suy tuyến tính}:
$$\mathbf{v} = \mathbf{p}_1 + \frac{c - f_1}{f_2 - f_1}(\mathbf{p}_2 - \mathbf{p}_1)$$
(5) Nối đỉnh thành tam giác theo bảng tra; chuyển sang khối kế. Bằng đối xứng (xoay + phần bù), 256 trường hợp giảm xuống \textbf{15 topology duy nhất}. Đầu ra: mesh tam giác xấp xỉ iso-surface.

\textbf{Câu 19. Volume rendering (DVR) vs Surface rendering.}\\
\textbf{Direct Volume Rendering} ray-cast xuyên khối, tích lũy màu/độ mờ qua \textbf{transfer function}: $C_{out} = \sum_i C_i \alpha_i \prod_{j<i}(1-\alpha_j)$. Hiển thị \textbf{toàn bộ cấu trúc bên trong bán trong suốt}, giữ tất cả dữ liệu, lý tưởng cho mô mềm và khám phá, nhưng chậm và khó đo lường. \textbf{Surface rendering (Marching Cubes)} trích mesh polygon ở một iso-value; render nhanh trên GPU, dễ đo/in 3D/cắt, nhưng \textbf{cố định một ngưỡng} và bỏ phần dữ liệu khác. \textbf{Dùng DVR} cho khám phá chẩn đoán, hiển thị đa mô; \textbf{dùng surface} cho lập kế hoạch phẫu thuật, in 3D, mô phỏng, phân tích trên mesh.

\textbf{Câu 20. Transfer function trong volume rendering là gì?}\\
\textbf{Transfer function} ánh xạ giá trị voxel vô hướng (và tùy chọn độ lớn gradient) sang \textbf{màu và độ mờ}: $\mathrm{TF}(v) \to (R,G,B,\alpha)$. Nó kiểm soát mô nào hiện ra: ví dụ gán $\alpha=0$ cho mô mềm và $\alpha=0.8$ đỏ cho xương trên CT để chỉ thấy bộ xương. \textbf{TF 1D} chỉ dùng cường độ; \textbf{TF 2D} thêm độ lớn gradient để làm nổi biên. Thiết kế TF tốt là \textbf{phần khó nhất của DVR} --- thường tương tác với widget histogram. Preset TF cho CT (xương, phổi, mạch) và MRI có sẵn trong viewer. TF về cơ bản là segmentation thị giác không cần nhãn rõ ràng.

\textbf{Câu 21. Iterative (MLEM, OSEM) vs FBP.}\\
\textbf{FBP} nhanh, giải tích, chính xác cho dữ liệu đầy đủ không nhiễu nhưng \textbf{khuếch đại nhiễu} (do ramp filter) và sinh streak khi liều thấp / ít góc. \textbf{MLEM} (Maximum Likelihood Expectation Maximization) mô hình thống kê photon (Poisson) và cập nhật lặp: $f^{(n+1)} = f^{(n)} \cdot A^T(p / Af^{(n)}) / A^T \mathbf{1}$. \textbf{OSEM} (Ordered Subset EM) tăng tốc MLEM bằng phân nhóm projection, tăng tốc $\sim$10 lần. Phương pháp lặp tích hợp \textbf{mô hình hệ thống} (hình học, tán xạ, suy giảm) và \textbf{regularization} (TV, bảo cạnh), cho SNR tốt hơn ở liều thấp, nhưng \textbf{chậm hơn} và có thể quá mịn. CT/PET hiện đại dùng iterative làm mặc định; tái tạo deep-learning đang nổi.

\textbf{Câu 22. SENSE và GRAPPA cho MRI song song.}\\
Parallel imaging tăng tốc MRI bằng lấy mẫu thưa k-space và dùng \textbf{mảng đa cuộn} (8--64 cuộn), mỗi cuộn có sensitivity không gian biết trước. \textbf{SENSE} (SENSitivity Encoding) hoạt động trong miền ảnh: k-space thưa cho ảnh aliased mỗi cuộn; SENSE giải hệ tuyến tính với sensitivity $C_i$ để tách aliasing tại mỗi pixel. \textbf{GRAPPA} (GeneRalized Autocalibrating Partial Parallel Acquisition) hoạt động trong k-space: dùng vùng \textbf{auto-calibration} đầy đủ để học trọng số tuyến tính tổng hợp các dòng k-space thiếu từ dòng thu kề trên các cuộn. Tăng tốc $R = 2$--$4$ là thường quy, giảm thời gian quét đổi lấy SNR giảm $\sqrt{R}$.

\textbf{Câu 23. Compressed sensing trong MRI là gì?}\\
\textbf{Compressed sensing (CS)} tái tạo ảnh MR từ \textbf{ít mẫu k-space hơn nhiều} so với Nyquist yêu cầu, khai thác hai sự thật: (1) ảnh MR \textbf{thưa} trong miền biến đổi nào đó (wavelet, total variation, sai phân hữu hạn); (2) lấy mẫu ngẫu nhiên/không tương quan trải aliasing thành nhiễu. CS giải:
$$\min_x \|\Psi x\|_1 \quad \text{s.t.} \quad \|F_u x - y\|_2 \leq \epsilon$$
với $\Psi$ là biến đổi sparsifying, $F_u$ Fourier thưa, $y$ k-space đo. Solver lặp phi tuyến (FISTA, ADMM) khôi phục $x$. CS cho phép tăng tốc 4--10 lần, đặc biệt cho MRI động (tim, cản quang). Kết hợp với parallel imaging và deep learning (variational network), nó nền tảng cho MRI nhanh hiện đại.

\textbf{Câu 24. Pipeline VTK cho hiển thị y học.}\\
\textbf{VTK} (Visualization Toolkit) dùng pipeline luồng dữ liệu: \textbf{Source $\to$ Filter $\to$ Mapper $\to$ Actor $\to$ Renderer $\to$ RenderWindow}. Với CT: (1) \texttt{vtkDICOMImageReader} nạp khối. (2) Filter tiền xử lý: \texttt{vtkImageThreshold}, \texttt{vtkImageGaussianSmooth}. (3) \texttt{vtkMarchingCubes} hoặc \texttt{vtkContourFilter} trích iso-surface; \textit{hoặc} \texttt{vtkSmartVolumeMapper} ray-cast DVR trên GPU. (4) Tùy chọn \texttt{vtkDecimatePro}, \texttt{vtkSmoothPolyDataFilter}. (5) \texttt{vtkPolyDataMapper} chuyển sang nguyên thủy đồ họa. (6) \texttt{vtkActor} thêm vật liệu/màu; \texttt{vtkVolume} cho DVR với \texttt{vtkColorTransferFunction} và \texttt{vtkPiecewiseFunction}. (7) \texttt{vtkRenderer} tổng hợp; \texttt{vtkRenderWindow} hiển thị; \texttt{vtkRenderWindowInteractor} xử lý input. 3D Slicer và ParaView xây trên VTK.

\textbf{Câu 25. Từ CT sang mesh xương in 3D --- pipeline.}\\
(1) \textbf{Thu} stack CT DICOM. (2) \textbf{Nạp} vào 3D Slicer / ITK-SNAP. (3) \textbf{Window/level} sang xương (HU $> +200$). (4) \textbf{Segment} xương bằng ngưỡng (HU $> 200$--300) cộng \textbf{region-growing} hoặc công cụ AI để loại mạch cản quang và bàn. (5) \textbf{Dọn hình thái}: closing, lấp lỗ, loại đảo. (6) \textbf{Marching Cubes} sinh mesh tam giác ở iso-value. (7) \textbf{Smooth} (Laplacian hoặc Taubin) và \textbf{decimate} về số tam giác hợp lý. (8) \textbf{Sửa}: đảm bảo manifold, watertight, không tự cắt (MeshLab, Meshmixer). (9) \textbf{Xuất STL/3MF}. (10) \textbf{Slice} cho in (Cura/PrusaSlicer) và in trên FDM/SLA. Dùng cho lập kế hoạch implant chỉnh hình, giảng dạy, surgical guide.

\textbf{Câu 26. Marching Cubes hỗ trợ lập kế hoạch phẫu thuật.}\\
Sau segmentation CT/MRI, \textbf{Marching Cubes} chuyển mỗi cơ quan có nhãn (u, mạch, xương) thành \textbf{mesh tam giác kín}. Bác sĩ có thể: (1) xem mesh 3D từ mọi góc, xoay thời gian thực trên GPU; (2) đo khoảng cách, thể tích, góc cho lập kế hoạch cắt bỏ; (3) in mô hình vật lý để diễn tập xúc giác; (4) đưa mesh vào \textbf{hệ VR/AR} chồng giải phẫu lên bệnh nhân; (5) lập quỹ đạo dụng cụ tránh cấu trúc quan trọng. So với hiển thị khối thô, mesh bề mặt \textbf{nhẹ, dễ thao tác, sẵn sàng in 3D}, biến MC thành công cụ chủ lực của lập kế hoạch phẫu thuật cá nhân hóa.

\textbf{Câu 27. Volume rendering trong lâm sàng.}\\
\textbf{Direct Volume Rendering} dùng cho: (1) \textbf{CT angiography} --- hiển thị mạch và hẹp 3D không cần segmentation; (2) \textbf{Nội soi đại tràng/phế quản ảo} --- bay xuyên tạng rỗng; (3) \textbf{Lập kế hoạch hàm mặt / chỉnh hình} --- tổng quan sọ và xương; (4) \textbf{Lập kế hoạch xạ trị} --- hiển thị target và OAR; (5) \textbf{Giao tiếp với bệnh nhân} --- ảnh 3D trực quan để đồng thuận. Ray-cast GPU hiện đại render khối 512$^3$ ở $>$30 fps. Cinematic rendering (path tracing, Monte Carlo) tăng độ chân thực. DVR bổ sung cho lát MPR bằng cách cho ngữ cảnh không gian mà không cần segmentation hoàn hảo.

\textbf{Câu 28. Tại sao tái tạo thời gian thực quan trọng trong can thiệp?}\\
\textbf{Can thiệp X-quang} (đưa ống thông, sinh thiết, đốt, phẫu thuật dưới fluoroscopy hoặc MR) cần ảnh \textbf{cập nhật trong mili giây} để bác sĩ phản ứng. Cụ thể: (1) \textbf{theo dõi dụng cụ} cần độ trễ $<$ 100 ms; (2) \textbf{chuyển động hô hấp/tim} đòi tốc độ khung $>$ 5 Hz; (3) \textbf{điều khiển vòng kín} (siêu âm hội tụ dẫn MR, phẫu thuật robot) cần độ trễ dưới khung xác định. Tái tạo chậm nghĩa là dẫn đường lỗi thời, kéo dài thời gian xạ/gây mê, và tăng nguy cơ biến chứng. Vì thế \textbf{FBP tăng tốc GPU, parallel imaging tái tạo nhanh, compressed sensing thời gian thực, và tái tạo dựa AI} là hướng nghiên cứu quan trọng.

\textbf{Câu 29. Xử lý ảnh phát hiện u trên MRI.}\\
\textbf{Pipeline phát hiện u} điển hình: (1) \textbf{Thu đa sequence} (T1, T1+Gd, T2, FLAIR, DWI/ADC) cho tương phản bổ sung. (2) \textbf{Tiền xử lý}: hiệu chỉnh bias-field (N4), loại sọ, đăng ký các sequence, chuẩn hóa cường độ. (3) \textbf{Trích đặc trưng}: cường độ voxel qua các sequence, texture (GLCM, LBP), gradient, atlas vị trí. (4) \textbf{Segmentation} bằng ngưỡng, region growing, level set, hoặc mô hình \textbf{CNN/U-Net} hiện đại (cuộc thi BraTS). (5) \textbf{Hậu xử lý}: thành phần liên thông, tính thể tích, vẽ contour. (6) \textbf{Radiomics}: trích đặc trưng hình dạng và texture cho tiên lượng. Đầu ra: mặt nạ u, thể tích, và độ tương đồng với tổn thương mẫu, giúp bác sĩ X-quang phát hiện, phân loại, và theo dõi u theo thời gian.

\end{paracol}
